\documentclass[10pt,a4paper]{amsart}
\usepackage{amsmath,amssymb,amsthm,graphicx}
% Retained preamble for the Stacks Project chapter snapshot in this directory.
% The chapter's own preamble.tex is a copy of the shared Stacks Project preamble; the block
% below is sliced MECHANICALLY, line for line, from it. Nothing was added, removed or
% redefined; the wrapper class, the xr-hyper/hyperref wiring and the externaldocument list
% are omitted because the assembler supplies the document class.
%
% lines taken from preamble.tex: 3, 5, 9, 11, 12, 156, 157, 158, ... (51 lines total)
%
% package coverage restored from the chapter's own preamble (2026-10-13): xy, xy-2cell

\usepackage[all]{xy}
\xyoption{2cell}
\UseAllTwocells

\newenvironment{reference}{\comment}{\endcomment}
\newenvironment{slogan}{\comment}{\endcomment}
\newenvironment{history}{\comment}{\endcomment}
\theoremstyle{plain}
\newtheorem{theorem}[subsection]{Theorem}
\newtheorem{proposition}[subsection]{Proposition}
\newtheorem{lemma}[subsection]{Lemma}
\theoremstyle{definition}
\newtheorem{definition}[subsection]{Definition}
\newtheorem{example}[subsection]{Example}
\newtheorem{exercise}[subsection]{Exercise}
\newtheorem{situation}[subsection]{Situation}
\theoremstyle{remark}
\newtheorem{remark}[subsection]{Remark}
\newtheorem{remarks}[subsection]{Remarks}
\numberwithin{equation}{subsection}
\def\lim{\mathop{\mathrm{lim}}\nolimits}
\def\colim{\mathop{\mathrm{colim}}\nolimits}
\def\Spec{\mathop{\mathrm{Spec}}}
\def\Hom{\mathop{\mathrm{Hom}}\nolimits}
\def\Ext{\mathop{\mathrm{Ext}}\nolimits}
\def\SheafHom{\mathop{\mathcal{H}\!\mathit{om}}\nolimits}
\def\SheafExt{\mathop{\mathcal{E}\!\mathit{xt}}\nolimits}
\def\Sch{\mathit{Sch}}
\def\Mor{\mathop{\mathrm{Mor}}\nolimits}
\def\Ob{\mathop{\mathrm{Ob}}\nolimits}
\def\Sh{\mathop{\mathit{Sh}}\nolimits}
\def\NL{\mathop{N\!L}\nolimits}
\def\CH{\mathop{\mathrm{CH}}\nolimits}
\def\proetale{{pro\text{-}\acute{e}tale}}
\def\etale{{\acute{e}tale}}
\def\QCoh{\mathit{QCoh}}
\def\Ker{\mathop{\mathrm{Ker}}}
\def\Im{\mathop{\mathrm{Im}}}
\def\Coker{\mathop{\mathrm{Coker}}}
\def\Coim{\mathop{\mathrm{Coim}}}
\DeclareMathSymbol{\boxtimes}{\mathbin}{AMSa}{"02}
\def\QCohstack{\mathcal{QC}\!\mathit{oh}}
\def\Cohstack{\mathcal{C}\!\mathit{oh}}
\def\Spacesstack{\mathcal{S}\!\mathit{paces}}
\def\Quotfunctor{\mathrm{Quot}}
\def\Hilbfunctor{\mathrm{Hilb}}
\def\Curvesstack{\mathcal{C}\!\mathit{urves}}
\def\Polarizedstack{\mathcal{P}\!\mathit{olarized}}
\def\Complexesstack{\mathcal{C}\!\mathit{omplexes}}
\def\Pic{\mathop{\mathrm{Pic}}\nolimits}
\def\Picardstack{\mathcal{P}\!\mathit{ic}}
\def\Picardfunctor{\mathrm{Pic}}
\def\Deformationcategory{\mathcal{D}\!\mathit{ef}}
\usepackage{cleveref}
\crefname{originallemma}{original lemma}{original lemmas}
\Crefname{originallemma}{Original lemma}{Original lemmas}
\makeatletter
\newlabel{descent-lemma-properties-quasi-coherent}{{0GNB}{source}{}{sourcepointer-1}{}}
\newlabel{descent-lemma-properties-quasi-coherent@cref}{{[originallemma][0][]0GNB}{[1][0][]source}{}{}{}}
\makeatother
\crefname{originallemma}{original lemma}{original lemmas}
\Crefname{originallemma}{Original lemma}{Original lemmas}
\makeatletter
\newlabel{sites-modules-lemma-internal-hom-restriction}{{0E8H}{source}{}{sourcepointer-2}{}}
\newlabel{sites-modules-lemma-internal-hom-restriction@cref}{{[originallemma][0][]0E8H}{[1][0][]source}{}{}{}}
\makeatother
\crefname{originallemma}{original lemma}{original lemmas}
\Crefname{originallemma}{Original lemma}{Original lemmas}
\makeatletter
\newlabel{sites-modules-lemma-pullback-internal-hom}{{0GN2}{source}{}{sourcepointer-3}{}}
\newlabel{sites-modules-lemma-pullback-internal-hom@cref}{{[originallemma][0][]0GN2}{[1][0][]source}{}{}{}}
\makeatother
\crefname{originallemma}{original lemma}{original lemmas}
\Crefname{originallemma}{Original lemma}{Original lemmas}
\makeatletter
\newlabel{sites-modules-lemma-tensor-product-pullback}{{03EL}{source}{}{sourcepointer-4}{}}
\newlabel{sites-modules-lemma-tensor-product-pullback@cref}{{[originallemma][0][]03EL}{[1][0][]source}{}{}{}}
\makeatother
\title{Stacks, Sheaves on Algebraic Stacks: lemma-lqc-colimits}
\author{The Stacks Project Authors}
\date{stacks-project:stacks-sheaves.tex:local-snapshot-20261012; GNU FDL 1.2 or later}
\begin{document}
\maketitle
\noindent\textit{Editorial note.} Statement span is lines 1977-2006 of sources/stacks-algebraic-stacks-sheaves/pinned-original-source.tex; proof span is lines 2008-2106, the complete author proof (97 body lines).\par\medskip
\noindent\textit{Editorial note.} Navigation stubs: 4 cited-but-not-included labels are supplied as non-load-bearing pointers with their pinned tags where the snapshot tag index has them.\par\medskip
\noindent\textit{Editorial note.} extra\_packages is empty on purpose: the retained preamble already loads xy with options and a second declaration clashes in the assembler wrapper.\par\medskip
\noindent\textit{Editorial note.} Text is retained verbatim, including the author's own typos and placeholders.\par\medskip
\noindent\textit{Editorial note.} Queue index 2 of queue\_algebraic-stacks-sheaves.json (pb=4649, stmt 1977-2006, proof 2008-2106).\par\medskip
\noindent\textit{Allowed standard prerequisites (editorial).}\begin{itemize}
\item Explicit prerequisite boundary. The cited inputs are prerequisites; the argument in the retained span is not replaced by any cited result.
\end{itemize}\medskip
\ifdefined\hypertarget\hypertarget{sourcepointer-1}{}\fi
\noindent\textit{Source reference.} Original lemma 0GNB: sources/stacks-algebraic-stacks-sheaves/pinned-original-source.tex, source lines 2106--2106. Cited by the retained statement or proof but not part of the collected spans; supplied only as a navigation pointer. The argument does not depend on its content.\par\medskip
\ifdefined\hypertarget\hypertarget{sourcepointer-2}{}\fi
\noindent\textit{Source reference.} Original lemma 0E8H: sources/stacks-algebraic-stacks-sheaves/pinned-original-source.tex, source lines 2106--2106. Cited by the retained statement or proof but not part of the collected spans; supplied only as a navigation pointer. The argument does not depend on its content.\par\medskip
\ifdefined\hypertarget\hypertarget{sourcepointer-3}{}\fi
\noindent\textit{Source reference.} Original lemma 0GN2: sources/stacks-algebraic-stacks-sheaves/pinned-original-source.tex, source lines 2106--2106. Cited by the retained statement or proof but not part of the collected spans; supplied only as a navigation pointer. The argument does not depend on its content.\par\medskip
\ifdefined\hypertarget\hypertarget{sourcepointer-4}{}\fi
\noindent\textit{Source reference.} Original lemma 03EL: sources/stacks-algebraic-stacks-sheaves/pinned-original-source.tex, source lines 2106--2106. Cited by the retained statement or proof but not part of the collected spans; supplied only as a navigation pointer. The argument does not depend on its content.\par\medskip
\begin{lemma}
\label{lemma-lqc-colimits}
Let $p : \mathcal{X} \to (\Sch/S)_{fppf}$ be a category fibred in
groupoids.
\begin{enumerate}
\item The category $\textit{LQCoh}(\mathcal{O}_\mathcal{X})$
has colimits and they agree with colimits in the category
$\textit{Mod}(\mathcal{X}_\etale, \mathcal{O}_\mathcal{X})$.
\item The category $\textit{LQCoh}(\mathcal{O}_\mathcal{X})$
is abelian with kernels and cokernels computed in
$\textit{Mod}(\mathcal{X}_\etale, \mathcal{O}_\mathcal{X})$,
in other words the inclusion functor is exact.
\item Given a short exact sequence
$0 \to \mathcal{F}_1 \to \mathcal{F}_2 \to \mathcal{F}_3 \to 0$ of
$\textit{Mod}(\mathcal{X}_\etale, \mathcal{O}_\mathcal{X})$
if two out of three are locally quasi-coherent so is the third.
\item Given $\mathcal{F}, \mathcal{G}$ in
$\textit{LQCoh}(\mathcal{O}_\mathcal{X})$
the tensor product $\mathcal{F} \otimes_{\mathcal{O}_\mathcal{X}} \mathcal{G}$
in $\textit{Mod}(\mathcal{X}_\etale, \mathcal{O}_\mathcal{X})$
is an object of $\textit{LQCoh}(\mathcal{O}_\mathcal{X})$.
\item Given $\mathcal{F}, \mathcal{G}$ in
$\textit{LQCoh}(\mathcal{O}_\mathcal{X})$
with $\mathcal{F}$ of finite presentation on
$\mathcal{X}_\etale$ the sheaf
$\SheafHom_{\mathcal{O}_\mathcal{X}}(\mathcal{F}, \mathcal{G})$
in $\textit{Mod}(\mathcal{X}_\etale, \mathcal{O}_\mathcal{X})$
is an object of $\textit{LQCoh}(\mathcal{O}_\mathcal{X})$.
\end{enumerate}
\end{lemma}

\begin{proof}
In the arguments below $x$ denotes an arbitrary object of $\mathcal{X}$
lying over the scheme $U$. To show that an object $\mathcal{H}$
of $\textit{Mod}(\mathcal{X}_\etale, \mathcal{O}_\mathcal{X})$ is
in $\textit{LQCoh}(\mathcal{O}_\mathcal{X})$ we will show that the
restriction $x^*\mathcal{H}|_{U_\etale} = \mathcal{H}|_{U_\etale}$
is a quasi-coherent object of $\textit{Mod}(U_\etale, \mathcal{O}_U)$.

\medskip\noindent
Proof of (1). Let $\mathcal{I} \to \textit{LQCoh}(\mathcal{O}_\mathcal{X})$,
$i \mapsto \mathcal{F}_i$ be a diagram.
Consider the object $\mathcal{F} = \colim_i \mathcal{F}_i$ of
$\textit{Mod}(\mathcal{X}_\etale, \mathcal{O}_\mathcal{X})$.
The pullback functor $x^*$ commutes with all colimits
as it is a left adjoint. Hence
$x^*\mathcal{F} = \colim_i x^*\mathcal{F}_i$. Similarly we have
$x^*\mathcal{F}|_{U_\etale} =
\colim_i x^*\mathcal{F}_i|_{U_\etale}$.
Now by assumption each $x^*\mathcal{F}_i|_{U_\etale}$
is quasi-coherent. Hence $\colim_i x^*\mathcal{F}_i|_{U_\etale}$
is quasi-coherent by
Descent, Lemma \ref{descent-lemma-properties-quasi-coherent}.
Thus $x^*\mathcal{F}|_{U_\etale}$ is quasi-coherent as desired.

\medskip\noindent
Proof of (2). It follows from (1) that cokernels exist in
$\textit{LQCoh}(\mathcal{O}_\mathcal{X})$ and agree with the cokernels computed
in $\textit{Mod}(\mathcal{X}_\etale, \mathcal{O}_\mathcal{X})$.
Let $\varphi : \mathcal{F} \to \mathcal{G}$ be a morphism of
$\textit{LQCoh}(\mathcal{O}_\mathcal{X})$ and let
$\mathcal{K} = \Ker(\varphi)$ computed in
$\textit{Mod}(\mathcal{X}_\etale, \mathcal{O}_\mathcal{X})$.
If we can show that $\mathcal{K}$ is a locally quasi-coherent module,
then the proof of (2) is complete. To see this, note that kernels
are computed in the category of presheaves (no sheafification necessary).
Hence $\mathcal{K}|_{U_\etale}$ is the kernel of the map
$\mathcal{F}|_{U_\etale} \to \mathcal{G}|_{U_\etale}$,
i.e., is the kernel of a map of quasi-coherent sheaves on $U_\etale$
whence quasi-coherent by
Descent, Lemma \ref{descent-lemma-properties-quasi-coherent}.
This proves (2).

\medskip\noindent
Proof of (3). Let
$0 \to \mathcal{F}_1 \to \mathcal{F}_2 \to \mathcal{F}_3 \to 0$
be a short exact sequence of
$\textit{Mod}(\mathcal{X}_\etale, \mathcal{O}_\mathcal{X})$.
Since we are using the \'etale topology, the restriction
$0 \to \mathcal{F}_1|_{U_\etale} \to
\mathcal{F}_2|_{U_\etale} \to \mathcal{F}_3|_{U_\etale} \to 0$
is a short exact sequence too. Hence (3) follows from the corresponding
statement in
Descent, Lemma \ref{descent-lemma-properties-quasi-coherent}.

\medskip\noindent
Proof of (4). Let $\mathcal{F}$ and $\mathcal{G}$ be in
$\textit{LQCoh}(\mathcal{O}_\mathcal{X})$. Since restriction to $U_\etale$
is given by pullback along the morphism of ringed topoi
$U_\etale \to (\Sch/U)_\etale \to \mathcal{X}_\etale$
we see that the restriction of the tensor product
$\mathcal{F} \otimes_{\mathcal{O}_\mathcal{X}} \mathcal{G}$
to $U_\etale$ is equal to
$\mathcal{F}|_{U_\etale} \otimes_{\mathcal{O}_U} \mathcal{G}|_{U_\etale}$,
see Modules on Sites, Lemma \ref{sites-modules-lemma-tensor-product-pullback}.
Since $\mathcal{F}|_{U_\etale}$ and $\mathcal{G}|_{U_\etale}$
are quasi-coherent, so is their tensor product, see
Descent, Lemma \ref{descent-lemma-properties-quasi-coherent}.

\medskip\noindent
Proof of (5). Let $\mathcal{F}$ and $\mathcal{G}$ be in
$\textit{LQCoh}(\mathcal{O}_\mathcal{X})$ with $\mathcal{F}$
of finite presentation. Since
$(\Sch/U)_\etale = \mathcal{X}_\etale/x$
is a localization of $\mathcal{X}_\etale$ at an object
we see that the restriction of
$\SheafHom_{\mathcal{O}_\mathcal{X}}(\mathcal{F}, \mathcal{G})$
to $(\Sch/U)_\etale$ is equal to
$$
\mathcal{H} =
\SheafHom_{\mathcal{O}|_{(\Sch/U)_\etale}}(
\mathcal{F}|_{(\Sch/U)_\etale}, \mathcal{G}|_{(\Sch/U)_\etale})
$$
by Modules on Sites, Lemma \ref{sites-modules-lemma-internal-hom-restriction}.
The morphism of ringed topoi
$(U_\etale, \mathcal{O}_U) \to ((\Sch/U)_\etale, \mathcal{O})$
is flat as the pullback of $\mathcal{O}$ is $\mathcal{O}_U$.
Hence the pullback of $\mathcal{H}$ by this morphism is equal to
$\SheafHom_{\mathcal{O}_U}(\mathcal{F}|_{U_\etale}, \mathcal{G}|_{U_\etale})$
by Modules on Sites, Lemma \ref{sites-modules-lemma-pullback-internal-hom}.
In other words, the restriction of
$\SheafHom_{\mathcal{O}_\mathcal{X}}(\mathcal{F}, \mathcal{G})$
to $U_\etale$ is
$\SheafHom_{\mathcal{O}_U}(\mathcal{F}|_{U_\etale}, \mathcal{G}|_{U_\etale})$.
Since $\mathcal{F}|_{U_\etale}$ and $\mathcal{G}|_{U_\etale}$
are quasi-coherent, so is
$\SheafHom_{\mathcal{O}_U}(\mathcal{F}|_{U_\etale}, \mathcal{G}|_{U_\etale})$,
see Descent, Lemma \ref{descent-lemma-properties-quasi-coherent}.
We conclude as before.
\end{proof}

\end{document}
